\subsection{Tensor products of Macaulay or Gorenstein algebras}

PROPOSITION 4. Let k be a field, A an essentially finitely generated k-algebra, and B a k-algebra.

\begin{enumerate}
\def\labelenumi{\alph{enumi})}
\item
  If A and B are Macaulay rings, then so is A ⊗k B.
\item
  If A and B are Gorenstein rings, then so is \(\Lambda\otimes_{k}\mathrm{B}\) .
\end{enumerate}

Suppose that A and B are Macaulay (resp. Gorenstein) rings, and let us prove that the same is true of \(A \otimes_{k} B\) . The ring \(A \otimes_{k} B\) is noetherian (no. 1, cor. of prop. 2). The A-module \(A \otimes_{k} B\) is free, hence flat. By prop. 10 of § 2, no. 7 (resp. cor. 1 of prop. 12 of § 3, no. 8), it suffices for us to prove that \(\kappa(\mathfrak{p}) \otimes_{k} B\) is a Macaulay ring (resp. Gorenstein ring) for every prime ideal p of A. The extension \(\kappa(\mathfrak{p})\) Since k is of finite type (no. 1, prop. 2 and example 1), we are therefore reduced to proving the statement in the case where the k-algebra A is a finitely generated extension K of k.

Let \((t_{1},\ldots,t_{n})\) a transcendence basis of K over k, so that K is a finite degree extension of the pure extension \(k^{\prime}=k(t_{1},\ldots,t_{n})\) (A, V, p.~112, prop. 17). The ring \(B^{\prime}=k^{\prime}\otimes_{k}B\) is isomorphic to a ring of fractions of the polynomial ring \(B[T_{1},\ldots,T_{n}]\) therefore is a Macaulay (resp. Gorenstein) ring by cor. 2 of prop. 10 of § 2, no. 7 (resp. cor. 3 of prop. 12 of § 3, no. 8). Since \(K\otimes_{k}B\) is identified with \(K\otimes_{k^{\prime}}B^{\prime}\) we are reduced to proving assertions a) and b) when K is a finite-degree extension of k.

The ring K ⊗k B is then a free B-module of finite rank, hence a Macaulay B-module. By § 2, no. 6, cor. 1 of prop. 8, it is a Macaulay ring, which proves a). Suppose now that B is a Gorenstein ring. There exist subextensions K \(_{i}\) with 0 ≤ i ≤ m, of K

\[
k = \mathrm{K} _ {0} \subset \mathrm{K} _ {1} \subset \dots \subset \mathrm{K} _ {m} = \mathrm{K}
\]

such that \(K_{i}\) let a \(K_{i-1}\) -monogenic algebra for \(i = 1, \ldots, m\) , and this allows us to reduce to the case where the extension K of k is monogenic (of finite degree). The canonical homomorphism \(B \to K \otimes_{k} B\) makes \(K \otimes_{k} B\) a flat B-algebra and, for every \(\mathfrak{q} \in \operatorname{Spec}(B)\) , the ring \((\mathrm{K} \otimes_{k} \mathrm{B}) \otimes_{\mathrm{B}} \mathfrak{k}(\mathfrak{q})\) , isomorphic to \(K \otimes_{k} \mathfrak{k}(\mathfrak{q})\) , is a \(\mathfrak{k}(\mathfrak{q})\) -monogenic algebra, hence a Gorenstein ring (§ 3, n° 9, example 5). Thus \(K \otimes_{k} B\) is a Gorenstein ring (§ 3, n° 8, cor. 1 of prop. 12), which completes the proof of b).

COROLLARY 1.--- Let \(k\) a field, K an extension of \(k\) and \(\Lambda\) an \(k\) -algebra essentially of finite type. For \(\mathrm{A}_{(\mathrm{K})}\) to be a Macaulay ring (resp. Gorenstein ring), it is necessary and sufficient that the same hold for A.

If A is a Macaulay (resp. Gorenstein) ring, then so is \(\mathrm{A}_{(\mathrm{K})}\) By prop. 4. Since \(\mathrm{A}_{(\mathrm{K})}\) is a faithfully flat A-module, the converse follows from prop. 10 of § 2, no. 7 (resp. from cor. 1 of prop. 12 of § 3, no. 8).

COROLLARY 2.--- Let k be a Noetherian ring, and let A and B be k-algebras. Assume that A is flat and essentially of finite type over k. If A and B are Macaulay rings (resp. Gorenstein rings), then A ⊗k B is a Macaulay ring (resp. Gorenstein ring).

Suppose first that B is a field; denote \(\varphi\) the canonical homomorphism from \(k\) into B, and \(\mathfrak{r}\) its kernel. The homomorphism \(\varphi\) induces a homomorphism from the field of fractions \(\kappa(\mathfrak{r})\) of \(k/\mathfrak{r}\) into B. Then A \(\otimes_{k}\) B is identified with (A \(\otimes_{k}\) K(r)) \(\otimes_{\kappa(\mathfrak{r})}\) B; as \(\varphi^{-1}(0)=\mathfrak{r}\) the ring A \(\otimes_{k}\) K(r) is a Macaulay (resp. Gorenstein) ring by prop. 10 of § 2, no. 7 (resp. cor. 1 of prop. 12 of § 3, no. 8). The assertion follows in this case from cor. 1.

Let us pass to the general case. The B-algebra A⊗ \(_{k}\) B is flat, and noetherian by the cor. of prop. 2 of no. 1. For each prime ideal q of B, the ring (A⊗ \(_{k}\) B)⊗ \(_{B}\) κ(q), which is identified with A⊗ \(_{k}\) κ(q) is a Macaulay (resp. Gorenstein) ring by the preceding results. One concludes by applying prop. 10 of § 2, no. 7 (resp. cor. 1 of prop. 12 of § 3, no. 8).

